%! Carrying Out a Test for the Difference of Two Population Proportions — Unit 6, Lesson 6.11 — Homework (Answer Key)
\documentclass[10pt]{article}
\usepackage{apstats-article}
\usepackage{apstats-key}   % pulls in -boxes; do NOT also load -boxes
\newcommand{\TallMath}[1]{$\displaystyle #1\rule[-1.4em]{0pt}{3.2em}$}

\begin{document}

\pageheader{Unit 6, Lesson 6.11}{Homework --- Answer Key}
\namedateperiod

\vspace{0.15in}

\begin{enumerate}

\item \textbf{Social Media and Academic Performance.}
High use: 54/180. Low use: 36/200. $\alpha=0.05$.

\begin{enumerate}[label=\alph*.]
  \item \ansline{$p_H$ = true proportion of high-use students with low GPA; $p_L$ = same for low-use.\\
        $H_0: p_H=p_L$; $H_a: p_H>p_L$ (one-sided right).}
  \item $\hat p_H=54/180=0.300$; $\hat p_L=36/200=0.180$; $\hat p_c=(54+36)/(180+200)=90/380\approx0.237$.\\
        Large Counts: $180(0.237)=42.7$, $180(0.763)=137.3$, $200(0.237)=47.4$, $200(0.763)=152.6$ --- \ans{all $\ge10$. \checkmark}\\
        Random/10\%: \ansline{Assumed both independently randomly sampled.}
  \item $SE=\sqrt{0.237(0.763)(1/180+1/200)}=\sqrt{0.18089 \times 0.01056}=\sqrt{0.001911}\approx0.04372$\\
        $z=(0.300-0.180)/0.04372=0.120/0.04372\approx\ans{2.74}$\\
        $p\text{-value}=P(Z>2.74)\approx\ans{0.003}$
  \item \ansline{Since $p$-value $= 0.003 \le \alpha = 0.05$, we reject $H_0$. We have convincing evidence that students who use social media more than 2 hours per day have a higher proportion of low-GPA outcomes than those who use it 2 hours or less per day.}
\end{enumerate}

\item \textbf{Employee Turnover.}
New: 18/240. Old: 30/240. $\alpha=0.05$.

\begin{enumerate}[label=\alph*.]
  \item \ansline{$p_N$ = true proportion of new-program employees who leave within a year; $p_O$ = same for old-program.\\
        $H_0: p_N=p_O$; $H_a: p_N<p_O$ (one-sided left: new program has lower turnover).}
  \item $\hat p_N=18/240=0.075$; $\hat p_O=30/240=0.125$; $\hat p_c=(18+30)/(480)=48/480=0.100$.\\
        Large Counts: $240(0.1)=24$, $240(0.9)=216$ --- \ans{all $\ge10$ for both groups. \checkmark}\\
        Random/10\%: \ansline{Assumed both samples $\le10\%$ of all employees in each program.}
  \item $SE=\sqrt{0.100(0.900)(1/240+1/240)}=\sqrt{0.090 \times 0.008\overline{3}}=\sqrt{0.000750}\approx0.02739$\\
        $z=(0.075-0.125)/0.02739=-0.050/0.02739\approx\ans{-1.83}$\\
        $p\text{-value}=P(Z<-1.83)\approx\ans{0.034}$
  \item \ansline{Since $p$-value $= 0.034 \le \alpha = 0.05$, we reject $H_0$. We have convincing evidence that the new onboarding program reduced employee turnover within the first year compared to the old program.}
\end{enumerate}

\end{enumerate}

\begin{reflectionbox}
\textbf{Unit 6 Complete.} You can now construct and interpret CIs and carry out significance
tests for one-sample and two-sample proportions. Great work!
\end{reflectionbox}

\end{document}
